\section{Discussion and Limitations}
\label{sec:discussion}

\method separates three sources of docking performance: candidate generation,
within-set selection, and molecular validity. High oracle-80 coverage indicates
that the fragment-level flow often reaches the native basin, whereas the
smaller top-1 rates quantify a ranking gap. The PoseBusters validity result
shows a second, partly independent gap that cannot be inferred from RMSD alone.
Registered selector and guidance studies did not pass their deployment gates,
which argues against attributing the remaining gap to a single post-processing
choice.

The rigid-fragment state reduces the dimensionality of ligand motion and
preserves within-fragment geometry, but independent fragment motions can still
distort cut-bond geometry or create protein clashes. The auxiliary atom and
distance-geometry losses mitigate these effects without enforcing a complete
molecular mechanics model. Large and highly flexible ligands remain the
clearest observed failure regime. Their lower oracle coverage shows that
improving only the confidence model will not resolve this weakness.

The pocket study further distinguishes receptor context from pocket
localization. Crops below \(8\,\angstrom\) remove useful receptor information,
whereas modest differences among \(8\), \(10\), and \(12\,\angstrom\) depend
on the dataset. In contrast, a \(2\,\angstrom\) center perturbation sharply
reduces both top-1 and oracle success. Robust deployment therefore requires an
independently validated pocket definition; random crop augmentation during
training does not make the model insensitive to localization error.

The present evaluation is reference-pocket redocking with a rigid receptor.
It does not test blind site identification, apo-to-holo receptor changes,
induced fit, affinity prediction, virtual screening, covalent docking, or
prospective generalization. In addition, the currently preserved training
split is a compatibility asset rather than the final strict external-exclusion
split. These limitations constrain the claims that can be made from the
current results. Literature values obtained with different pocket inputs,
candidate budgets, receptor preparation, or post-processing are therefore not
used as matched baselines in the main table.

The highest-priority next evidence is a strict ligand- and pocket-disjoint
training split with frozen external exclusions, followed by matched reruns of
representative classical and learned baselines. Independent training or
sampling seeds, wall-clock and memory measurements, and component ablations
for fragmentation, Newton--Euler aggregation, and auxiliary geometry losses
are also required to isolate which design choices cause the observed behavior.
Target-independent pocket centers and apo or predicted receptors would extend
the evidence beyond the present holo redocking boundary.
